% © 2026 Ing. David Jaroš — CC BY-NC-ND 4.0
%
% This work is licensed under a Creative Commons Attribution-NonCommercial-NoDerivatives
% 4.0 International License (CC BY-NC-ND 4.0).
%
% See LICENSE.md for full license text.
%
% Classification: [STD] sections 1-3, [MC] section 4, [MC/SPECULATIVE] section 5.

\documentclass[12pt,a4paper]{article}
\usepackage[a4paper, margin=2.5cm]{geometry}
\usepackage{amsmath,amssymb}
\usepackage{hyperref}
\title{\textbf{Kanonická kvantizace gravitonových módů v UBT}}
\author{Ing.\ David Jaro\v{s}}
\date{2026-05-11}

\begin{document}
\maketitle

\section{Linearizace $S[\Theta]$ okolo Schwarzschildova pozadí \,[STD]}

Vycházíme z kanonické akce
\[
S[\Theta]=\int d^4x\sqrt{-g}\left[\frac{R}{2\kappa}+\mathrm{Tr}\!\left((D_\mu\Theta)^\dagger(D^\mu\Theta)\right)-V(\Theta)\right].
\]
Položíme
\[
\Theta=\Theta_{\mathrm{Schw}}+\delta\Theta
\]
a rozvoj do druhého řádu dává kvadratickou akci $S_2[\delta\Theta]$, která
popisuje gravitonové fluktuace nad Schwarzschildovým pozadím.

\section{Identifikace módů \,[STD]}

Použijeme uzávěr z
\texttt{canonical/gr\_closure/zerilli\_derivation.tex}:
\begin{itemize}
  \item sudá parita $\delta\Theta_{\mathrm{even}} \to V_Z(r)$ (Zerilli),
  \item lichá parita $\delta\Theta_{\mathrm{odd}} \to V_{RW}(r)$ (Regge--Wheeler).
\end{itemize}
Po rozkladu do sférických harmonik zavádíme kanonické souřadnice
$Q_{\ell m}(t)$ a momenta $P_{\ell m}=\dot Q_{\ell m}$.

\section{Kanonická kvantizace \,[STD]}

Postulujeme komutační relace
\[
[Q_{\ell m}(t),P_{\ell' m'}(t)]=i\hbar\,\delta_{\ell\ell'}\delta_{mm'}.
\]
Pro každý mód:
\[
H_{\ell m}=\frac12\!\left(P_{\ell m}^2+\omega_{\ell m}^2Q_{\ell m}^2\right),
\]
kde $\omega_{\ell m}$ je určeno potenciálem $V_Z$ nebo $V_{RW}$.
Výsledkem je Fockova báze $|n_{\ell m}\rangle$, vakuum $|0\rangle$ a nulový
příspěvek $\frac12\hbar\omega_{\ell m}$ na mód.

\section{UBT specifika: společné pole $\Theta$ \,[MC]}

V UBT gravitační i gaugeové excitace vznikají ze stejného pole $\Theta$,
tedy
\[
\mathcal H_{\mathrm{grav}}\subset \mathcal H_\Theta.
\]
Pracovní podmínka pro další odvození:
\begin{itemize}
  \item generátory $\mathrm{SU}(2)_L\times\mathrm{U}(1)_Y$ působí na interní
  algebraickou část stavů,
  \item gravitonové módy působí na geometrickou (paritní/sférickou) část,
  \item přirozené super-výběrové sektory mohou vznikat z dekompozice
  algebry $\mathbb C\otimes\mathbb H$.
\end{itemize}
Explicitní komutační algebra mezi těmito sektory zůstává otevřená.

\section{Vakuová energie a kosmologická konstanta \,[MC/SPECULATIVE]}

Formální vakuová energie:
\[
E_{\mathrm{vac}}=\sum_{\ell,m}\frac12\hbar\omega_{\ell m},
\]
což je UV divergentní suma. Konzistentní kandidát regularizace je
$\zeta$-funkce, v návaznosti na postupy použité v alpha tracku
(\texttt{step2\_vacuum\_polarization.tex}).

Otevřený test:
\begin{itemize}
  \item zda $\zeta$-regularizovaná hodnota dává nenulový efektivní příspěvek
  do kosmologické konstanty v UBT,
  \item jak tento příspěvek souvisí s běžnou renormalizací gravitačního sektoru.
\end{itemize}

\end{document}

